\chapter{Построение дискретных командных генераторов}
\label{ch:chap3}

Синтезируем командный генератор гармонического сигнала вида 
\begin{equation}
	g(k) = A \sin (kT\omega),
\end{equation}

где $T=0.2$ с, $A=1$, $\omega =0.08$.

Дискретная модель внешнего воздействия описывается в пространстве состояний системой разностных уравнений

\begin{equation}
	\begin{cases}
		\xi (k+1) = \Gamma_d \xi (k),\\
		g(k) = H \xi (k),
	\end{cases}
\end{equation}

где $\xi(k)$ -- n-мерный вектор состояния дискретного командного генератора, $g(k)$ -- дискретное внешнее воздействие, $\Gamma$ -- матрица динамики дискретного генератора размерности $n \times n$, $H$ -- матрица выходов дискретного генератора размерности $1 \times n$.

Выберем в качестве первой координаты вектора состояния дискретного генератора сам сигнал:

\begin{equation}
	\xi_1(k) = g(k)
\end{equation}

Вычислим первую разность

\begin{multline}
	g(k+1) = A \sin ((k+1)T\omega) =  A \sin (kT\omega) \cos (T\omega) + A \cos(kT\omega) \sin(T\omega)=\\=
	g(k) \cos (T\omega) + A \cos(kT\omega) \sin(T\omega)
\end{multline}
сразу выразим

\begin{equation}
	 A \cos(kT\omega) = \frac{g(k+1)-g(k) \cos (T\omega)}{\sin(T\omega)}
\end{equation}

Перейдем ко второй разности

\begin{multline}
	g(k+2) = g(k+1) \cos (T \omega) + \sin (T \omega) \left(A \sin(kT\omega) \sin(T\omega)+A \cos(kT\omega) \cos(T\omega) \right)=\\=
	2g(k+1) \cos(T\omega) - g(k)
\end{multline}

Вектор состояния запишем как

\begin{equation}
	\begin{cases}
		\xi_1 (k) = g(k)\\
		\xi_2(k) = \xi_1(k+1) +g(k+1)\\
		\xi_3(k) = \xi_2(k+1) + g(k+2)
	\end{cases}
\end{equation}

Запишем модель в пространстве состояний

\begin{equation}
	\begin{cases}
		\begin{bmatrix}
			\xi_1 (k+1)\\
			\xi_2(k+1)
		\end{bmatrix} = 
		\begin{bmatrix}
			0 & 1\\
			-1 & 2 \cos(T\omega)
		\end{bmatrix}	\begin{bmatrix}
		\xi_1 (k)\\
		\xi_2(k)
		\end{bmatrix}\\
		g(k) = \begin{bmatrix}
			1 & 0
		\end{bmatrix}
			\begin{bmatrix}
			\xi_1 (k)\\
			\xi_2(k)
		\end{bmatrix}
	\end{cases}
\end{equation}

Начальные условия запишем в виде
\begin{equation}
	\begin{cases}
		\xi_1(0) = A \sin (0) = 0\\
		\xi_2 (0) = A \sin(0) \cos(T \omega) + A \cos(0) \sin(T\omega) = A \sin (T \omega)
	\end{cases}
\end{equation}

\endinput